\begin{figure}[ht]
\centering
\begin{tikzpicture}[
    lbl/.style={font=\scriptsize, text=black!55},
    dimlbl/.style={font=\tiny, text=black!45},
    clbl/.style={font=\scriptsize\bfseries},
]
\def\c{0.095}

% --- Helper: draw 21x16 face grid ---
% #1=x-offset, #2=bg (outside face), #3=face fill,
% #4=feature fill, #5=outline fill, #6=draw color
\newcommand{\drawface}[6]{%
    \foreach \x in {0,...,20} {
        \foreach \y in {0,...,15} {
            % Distance from face center (10, 7.5)
            \pgfmathsetmacro{\dist}{sqrt((\x-10)*(\x-10)+(\y-7.5)*(\y-7.5))}
            \pgfmathtruncatemacro{\inface}{\dist < 6.3 ? 1 : 0}
            \pgfmathtruncatemacro{\onedge}{(\dist >= 6.3)*(\dist < 7.2) ? 1 : 0}
            % Eyes (2x2) + smile curve
            \pgfmathtruncatemacro{\feat}{(
                (\x>=7)*(\x<=8)*(\y>=9)*(\y<=10) +
                (\x>=12)*(\x<=13)*(\y>=9)*(\y<=10) +
                (\x==6)*(\y==5) + (\x==14)*(\y==5) +
                (\x==7)*(\y==4) + (\x==13)*(\y==4) +
                (\x>=8)*(\x<=12)*(\y==3)
            ) > 0 ? 1 : 0}
            \ifnum\onedge=1
                \fill[#5, draw=#6]
                    ({#1+\x*\c},\y*\c) rectangle ++(\c,\c);
            \else\ifnum\inface=0
                \fill[#2, draw=#6]
                    ({#1+\x*\c},\y*\c) rectangle ++(\c,\c);
            \else\ifnum\feat=1
                \fill[#4, draw=#6]
                    ({#1+\x*\c},\y*\c) rectangle ++(\c,\c);
            \else
                \fill[#3, draw=#6]
                    ({#1+\x*\c},\y*\c) rectangle ++(\c,\c);
            \fi\fi\fi
        }
    }
}

% === Color image ===
\definecolor{imgbg}{rgb}{0.92, 0.92, 0.92}
\definecolor{ybg}{rgb}{0.95, 0.82, 0.12}
\definecolor{yfg}{rgb}{0.15, 0.10, 0.08}
\definecolor{yol}{rgb}{0.30, 0.25, 0.10}
\drawface{0}{imgbg}{ybg}{yfg}{yol}{black!4}
\node[lbl] at ({10.5*\c}, {16*\c + 0.22}) {color image};
\node[dimlbl, anchor=north] at ({10.5*\c}, -0.08)
    {$210 \times 160 \times 3$};

% === Arrow 1 ===
\draw[->, thick, black!35]
    ({21*\c + 0.12}, {8*\c}) -- ++(0.3, 0);

% === R channel ===
\def\rx{2.45}
\drawface{\rx}{red!6}{red!80}{red!12}{red!40}{red!8}
\node[clbl, text=red!70] at ({\rx + 10.5*\c}, {16*\c + 0.22}) {R};

% === G channel ===
\pgfmathsetmacro{\gx}{\rx + 21*\c + 0.06}
\drawface{\gx}{green!5}{green!70}{green!8}{green!35}{green!8!black!3}
\node[clbl, text=green!55!black]
    at ({\gx + 10.5*\c}, {16*\c + 0.22}) {G};

% === B channel ===
\pgfmathsetmacro{\bx}{\gx + 21*\c + 0.06}
\drawface{\bx}{blue!4}{blue!14}{blue!8}{blue!10}{blue!5}
\node[clbl, text=blue!60] at ({\bx + 10.5*\c}, {16*\c + 0.22}) {B};

% Brace under channel group
\draw[decorate, decoration={brace, amplitude=3pt, mirror}, black!35]
    (\rx, -0.05) -- ({\bx + 21*\c}, -0.05)
    node[midway, below=4pt, lbl] {3 channels};

% === Arrow 2 ===
\pgfmathsetmacro{\arrtwo}{\bx + 21*\c + 0.12}
\draw[->, thick, black!35]
    (\arrtwo, {8*\c}) -- ++(0.3, 0);

% === Grayscale output ===
\pgfmathsetmacro{\ox}{\arrtwo + 0.42}
\drawface{\ox}{black!6}{black!20}{black!85}{black!55}{black!5}
\node[lbl] at ({\ox + 10.5*\c}, {16*\c + 0.22}) {grayscale};
\node[dimlbl, anchor=north] at ({\ox + 10.5*\c}, -0.08)
    {$210 \times 160 \times 1$};

\end{tikzpicture}
\caption{Grayscale conversion. Three color channels combine into
  one brightness value per pixel, preserving spatial structure.}
\label{fig:grayscale}
\end{figure}
